\section{Introduction}\label{sec:introduction}

Somatic copy number alterations (CNAs), ranging from focal amplifications to gains and losses of whole chromosome arms, are among the most frequent genetic changes in cancer, and some are characteristic of particular tumour types. Examples are the co-deletion of chromosome arms 1p and 19q in oligodendroglioma and the gain of chromosome 7 with loss of chromosome 10 in glioblastoma \citep{tirosh2016oligo, patel2014}. Copy number is usually measured from bulk DNA, which averages over all cells in a sample. Measuring it in single cells adds two kinds of information. First, it identifies the malignant cells of a tumour, which is not always possible from expression alone, because malignant cells can closely resemble their non-malignant counterparts. Second, it resolves genetic subclones and relates them to the transcriptional states and microenvironment of the same cells.

CNAs were first inferred from single-cell transcriptomes in a study of primary glioblastoma \citep{patel2014}. A gain or loss of a large genomic region shifts the expression of many neighbouring genes in the same direction. The expression of any single gene reflects its copy number only weakly, but the average over many genes ordered along a chromosome reveals the change. In that study, log-transformed expression values were centred on each gene's mean, averaged in a moving window of 100 genes along the genome, and centred again within each cell. No normal reference cells had been isolated; the few normal cells in the data were recognised because their profiles clustered with an averaged profile of bulk normal brain. Their mean profile then served as the baseline, against which the malignant cells showed the gain of chromosome 7 and loss of chromosome 10 typical of glioblastoma. Applied to cancer cell lines with both RNA-seq and SNP array data, the same smoothing gave profiles that correlated with array-based copy number (r = 0.72) \citep{patel2014}.

The approach, implemented by Itay Tirosh in MATLAB, was refined in a series of single-cell tumour studies. In oligodendroglioma, relative expression values were capped at $\pm$3 before smoothing so that single genes could not dominate a window. Separate baselines were also defined for each normal cell type, and a value was counted as altered only if it lay outside the range of those baselines by more than a small margin \citep{tirosh2016oligo}. Studies of IDH-mutant glioma and head and neck cancer then classified cells as malignant by two scores: the overall magnitude of a cell's CNA profile, and its correlation with the average profile of the tumour \citep{venteicher2017, puram2017}. A later glioblastoma study defined subclones from chromosome arms with bimodal CNA values \citep{neftel2019}. The method was used to identify malignant cells in further studies of melanoma \citep{tirosh2016melanoma} and H3K27M-mutant glioma \citep{filbin2018}. It was later released as the R package infercna \citep{laffy_infercna} and used to verify malignant cells in a pan-cancer atlas \citep{gavish2023}.

As single-cell profiling of tumours spread, there was a need for an open-source tool with which any group could apply the method routinely. We developed inferCNV for this purpose, as an R package based on Tirosh's method, and distributed it through GitHub and later Bioconductor \citep{tickle2019infercnv, infercnvBioc}. inferCNV kept the core of the original procedure: the log-transformed, depth-normalised expression, the cap on relative values, smoothing over a window of about 100 genes, per-cell centring, and subtraction bounded by the range of several normal reference groups. It began as a visualisation tool, from which CNAs were read by inspecting heatmaps. It was then extended to predict discrete copy number states with a hidden Markov model (HMM), building on the HMM of HoneyBADGER \citep{fan2018}. The HMM's states are calibrated by simulating spiked-in CNAs that match the depth and noise of each data set \citep{haas2021thesis}. inferCNV also assesses its calls with a Bayesian mixture model and partitions cells into CNA-based subclusters. It does not, however, include the score-based classification of malignant cells used in the original studies.

inferCNV has since been applied to many tumour types and platforms. Early uses that named the package were in melanoma \citep{jerbyarnon2018} and synovial sarcoma \citep{jerbyarnon2021}. It has been used with data from Smart-seq2 \citep{maynard2020}, droplet-based 10x Genomics \citep{ma2019, bi2021, steen2021}, Seq-Well \citep{raghavan2021} and single-nucleus RNA-seq \citep{slyper2020, terekhanova2023}. Applications include studies of uveal melanoma with single-cell DNA validation \citep{durante2020}, high-grade serous ovarian cancer \citep{izar2020}, breast cancer \citep{wu2021breast} and small cell lung cancer \citep{chan2021}, and a pan-cancer analysis of recurrent cell states \citep{barkley2022}. A full-text search of Europe PMC on 20 October 2026 found 2,741 articles that mention ``inferCNV'', rising from 3 published in 2018 to 758 in 2025. These are mentions in titles, abstracts or full texts rather than citations, and an undercount, because full text can be searched only for the articles that Europe PMC holds in full. The method has also been carried into spatial transcriptomics. SpatialInferCNV applies inferCNV to spatially barcoded spots to reconstruct clonal CNAs in situ, including in histologically benign tissue \citep{erickson2022}, and has been used by others \citep{pei2025}. inferCNV has also been applied directly to spot-based data \citep{wu2021liver, denisenko2024, lucas2024, mo2024, harwood2024} and to spatially barcoded single-nucleus data \citep{russell2024}.

Several other methods infer CNAs from single-cell transcriptomes; the main differences between them are the use of a normal reference and of allelic information. Among methods that use expression alone, some require no reference. CONICSmat fits mixture models to the expression of each chromosome arm \citep{muller2018conics}. CopyKAT identifies diploid cells itself and segments the genome of each cell \citep{gao2021copykat}. SCEVAN classifies cells and segments clonal profiles \citep{defalco2023}. sciCNV \citep{mahdipourshirayeh2022scicnv} and the Python reimplementation infercnvpy \citep{sturm_infercnvpy} produce continuous profiles. Allele-aware methods add the allelic imbalance of heterozygous variants in the reads, which can show loss of heterozygosity even when expression is unchanged. HoneyBADGER combines an expression HMM with an allele model \citep{fan2018}, and CaSpER uses allele shifts to refine expression signals at several scales \citep{serinharmanci2020}. Numbat phases variants with a population haplotype panel and infers a phylogeny of subclones \citep{gao2023numbat}, and XClone infers states separately from read depth and from phased allele frequencies and then combines them \citep{huang2024xclone}. Methods for spatial data model neighbouring spots jointly, sometimes with allelic information \citep{elyanow2021starch, ma2024calicost}.

Independent benchmarks have compared inferCNV with these methods. \citet{schmid2025} compared six callers on 15 droplet-based human data sets with DNA-based truth. inferCNV (v1.10) was among the methods with the highest maximal F1, although differences between methods were mostly not significant. The authors found that it ``performed relatively well in all evaluated categories''. It also had the longest run time and the second-highest memory use, was more sensitive to gains than to losses, and depended strongly on a matched reference. Other benchmarks found it among the best methods for identifying tumour subpopulations within a platform but sensitive to batch and platform effects \citep{chen2025}. One found it reliable in separating tumour from normal cells, except at very high tumour purity when run without reference cells \citep{song2025}. A preprint comparing twelve methods found it accurate in classifying malignant cells but, together with one other method, the least scalable, failing to finish 10,000 cells within 24 hours, and found that the default Leiden subclustering over-partitioned cells \citep{chang2026preprint}. Despite this wide use, there has been no publication describing the algorithm and its implementation, and users have cited the GitHub repository or the Bioconductor package \citep{tickle2019infercnv, infercnvBioc}.

This article provides that description. We describe the algorithm and the inferCNV release in which it is implemented (Section~\ref{sec:implementation} and Supplementary Methods). We then illustrate its behaviour on simulated data with known CNAs and on tumour samples with DNA-based copy number measurements (Section~\ref{sec:results}).
